\begin{wrapfigure}[]{l}{.35\textwidth}
\vspace{-13pt}
\includegraphics[width=.35\textwidth]{../pictures/unicode-fool.jpeg}
\vspace{-28pt}
\end{wrapfigure}

\noindent Unless there is a new person to talk to, a lot of the ``education
stuff'' we do could grow stale. Many of the patterns and use cases for
peeragogy assume that there will be an audience or a new generation of
learners - hence the drive to create a guide. Note that the newcomer and
the \href{http://peeragogy.org/practice/wrapper/}{wrapper} may work
together to make the project accessible. Even in the absence of actual
newcomers, we're often try to look at things with a
``beginner's mind.'' 

\begin{quote}
\textbf{Régis Barondeau}: I joined this handbook project late, making me
a ``newcomer''. When I started to catch up, I rapidly faced doubts:
\emph{Where do I start? How can I help? How will I make it, having to
read more than 700 posts to catch up? What tools are we using ? How do I
use them?} Etc. Although this project is amazingly interesting, catching
the train while it already reached high speed can be an extreme sport.
By taking care of newcomers, we might avoid loosing valuable contribuors
because they don't know how and where to start, and keep our own project
on track.
\end{quote}

\begin{quote}
\textbf{Charlotte Pierce}: Joe was working a lot on the book,
and I thought ``this is interesting hard work, and he shouldn't have to
do this alone.'' As a Peeragogy newcomer, I was kindly welcomed and
mentored by Joe, Howard, Fabrizio, and others. I asked naive questions
and was met with patient answers, guiding questions, and resource links.
Concurrently, I bootstrapped myself into a position to contribute to the
workflow by editing the live manuscript for consistency, style, and
continuity.
\end{quote}
